\section{Discussion}
\label{sec:discussion}

\subsection{Answers to the Research Questions}
\emph{Evaluation protocol.} On these data the split protocol, not the model, decides whether scores look near-perfect: random row-level splits---with or without oversampling---give ROC-AUC above 0.99 because adjacent seconds of a session are near-duplicates with the same label, whereas session-grouped evaluation of the same columns gives \LeakGrpRocAuc. This is a known pitfall rather than a new finding, but it means that accuracy figures on gameplay cybersickness data should be read together with the split protocol.
\emph{RQ1.} On logged data the asymmetric reward encodes a cost-sensitive operating point: the greedy policy is the Bayes rule $P(y=1\mid s)\ge1/6$ of 5:1 cost-sensitive classification, and bootstrapping cannot improve it (Section~\ref{sec:formulation}).
\emph{RQ2.} Under leakage-controlled evaluation with the same search procedure and decision rules, E-DDQN does not outperform the standard DQN, has a higher cost than a cost-weighted random forest in every outer fold (significant under the group bootstrap, but not after Holm adjustment under the more conservative exact sign-flip test), and ranks seconds worse than elapsed time alone, which no model outperformed significantly; no model meets a safety-prioritized target.
\emph{RQ3.} As a sanity check, bounded manipulation of the telemetry channels barely changes threshold-free ranking, which is not evidence of robustness (Section~\ref{sec:attack}), while an adaptive adversary with label knowledge can still lower recall at fixed thresholds by \AttEvasionDropMin--\AttEvasionDropMax{} percentage points.

\subsection{Is E-DDQN Reinforcement Learning?}
With logged data the alarm cannot influence the subsequent telemetry, so the decision problem is a contextual bandit. E-DDQN remains a valid value-based learner of a cost-sensitive detection policy, but (i) the asymmetric reward is mathematically equivalent to a 5:1 cost-sensitive Bayes threshold, so it is a statement about the operating point rather than a different learning objective; (ii) bootstrapping cannot improve the optimal policy, and empirically removing it with E-DDQN's configuration held fixed changed little, whereas E-DDQN's under-performance is more consistent with noisy validation-based selection of configurations that use the questionnaire answers (Section~\ref{sec:why}); and (iii) under bandit feedback the agent observes the reward of one action only, although labels reveal both, which is why full-information replay is the principled way to use the labels. The same reasoning applies to any DQN-based classifier trained on a fixed dataset whose episodes end only with the data: its benefit has to come from estimation effects, and it should be compared with cost-weighted supervised models at the same operating point. Formulations that end an episode when a minority-class sample is misclassified~\cite{lin2020imbalanced,yang2024drlimbalanced} are action-dependent, but the termination acts only as an additional, $\gamma$-dependent penalty on that error (Section~\ref{sec:formulation}), so the same comparison is needed to show that they add more than a different operating point. Adaptive sequential control would require an intervention in the loop---for example an alarm that triggers a rendering mitigation whose effect on later telemetry and symptoms is measured---which no public cybersickness dataset provides.

\subsection{The Safety Objective}
A 5:1 cost makes an alarm the cheaper decision whenever $P(y=1\mid s)\ge1/6$. With a prevalence of 0.44 and weakly informative inputs, almost every second satisfies this, so the cost-optimal policy of every model is close to ``always alarm'': recall is high, but \DOneFPRMin--\DOneFPRMax\,\% of non-event seconds also trigger an alarm. High recall alone is therefore not evidence of a useful safety monitor. Judged on cost, specificity at matched recall, threshold-free ranking and cost curves across ratios, the intended safety objective---near-complete detection with a controlled false-alarm rate---is not achieved by any model on these data. The exposure-time baselines suggest where the obstacle lies: once sessions and linkage groups are not shared between training and test, no model extracted significantly more signal from game telemetry than how long the session has lasted, which distinguishes individual susceptibility only coarsely.

\subsection{The Adversarial Threat Model}
Rendering manipulations that induce sickness are documented~\cite{odeleye2021framerate,tseng2022dark}, and optimization-based input attacks can disable deep cybersickness detectors~\cite{kundu2025securing,ahmed2025adversarialvr}. UFFCSData, however, contains no manipulated sessions, so its label marks cybersickness, not an attack. Our stress tests therefore answer a narrower question than the cited attack studies: how far bounded, physically motivated changes of the logged frame-rate, orientation and locomotion channels move the detectors' decisions. The small effect on ranking is not evidence of a robust design: it reflects the inputs that the selected configurations use and the fact that a perturbation applied to every test session alike can shift scores without reordering them, while at fixed thresholds an adaptive adversary still changes decisions. A detector that used the telemetry more would also be easier to manipulate through it. Demonstrating detection of real rendering attacks requires data in which manipulations are applied to participants under ethical approval and their physiological effects are measured.

\subsection{Limitations}
\label{sec:limits}
(i) A single public dataset with \nSessionsEval{} curated sessions, \nGroupsEval{} independent groups and \nEventSessions{} onsets; confidence intervals are wide, fold-to-fold variation is large, and inference with so few clusters is approximate: the group bootstrap and the exact sign-flip test agree in direction but differ in strength (E-DDQN vs.\ random forest, Holm-adjusted $p=$\,\PHHolmRF{} vs.\ \PHSfRFHolm). (ii) Participant identity cannot be verified; linkage groups are conservative but cannot exclude a participant who answered the questionnaire differently in non-consecutive sessions, and they may merge different people with identical profiles. (iii) The label is a coarse self-report, logged at 1\,Hz but changing rarely within a session, whose elicitation procedure (prompted or spontaneous reports) is not documented, and the binary threshold (discomfort $\ge1$) merges slight with severe discomfort; a safety monitor might target grades 2--3 only, which occur in just \PHGradeFlightSevere{} Flight samples. (iv) The exposure-time baselines, the Holm adjustment, the grouping sensitivity analyses and the cost curves were specified after the primary results were known; they were committed before being run and use the same protocol, but they are post hoc. The within-fold ROC-AUC, the sign-flip tests and the feature-set analysis were added after a pre-submission review, when all other results were known. (v) Perturbations are synthetic and feature-space, their physiological effect is unknown, and the adversary is weaker than optimization-based attacks. (vi) Hyper-parameter search used a modest budget (\HPOBudget{} configurations per family; 12 for logistic regression) that covers a much smaller fraction of the RL search spaces (about \PHSpaceEDDQNMillion{} million configurations for E-DDQN) than of the random-forest space (\PHSpaceRF); a larger budget could change the ranking of the neural models, although the selection-optimism analysis suggests that additional search with the same noisy criterion would mostly increase optimism. (vii) The October-2019 sessions were used only for development and could not serve as an external test set because their camera-orientation channels are constant within every session (a logging defect). (viii) The diagnosis of E-DDQN's under-performance (Section~\ref{sec:why}) is exploratory: feature sets and discount factors were selected, not assigned, so no analysis isolates a cause.

\subsection{Implications}
For cybersickness detection from session data, results should be reported with participant- or session-grouped splits, alongside trivial policies and an exposure-time baseline, because here no model significantly outperformed elapsed time alone. For safety-prioritized detection, an asymmetric reward should be treated as what it is---a cost-sensitive operating point---and RL detectors should be compared with cost-weighted supervised models at matched recall or cost. Claims of robustness to manipulation should state which inputs the detector actually uses and evaluate decisions at fixed thresholds, since an unchanged ranking is not protection. Under grouped evaluation, participant-profile features deserve particular care, because with few groups they can fit group-level label rates that do not transfer.
